\section*{Chapter \ref{ch:b1:plate-kinematics} --- Plate Kinematics on a Sphere}

\begin{solution}{exo:b1:plate-kinematics:1}
(a) $\qty{0.210}{\degree\per\Myr} = 0.210\times\pi/180 = \qty{3.665e-3}{rad}$
per million years, that is \qty{3.665e-9}{rad.yr^{-1}}; with
$\qty{1}{yr} = \qty{3.156e7}{s}$, \qty{1.161e-16}{rad.s^{-1}}. (b) $360/0.210 =
\qty{1714}{\Myr}$, more than a third of the age of the Earth. (c) $\omega R =
\num{3.665e-9}\times\qty{6.371e9}{mm} = \qty{23.35}{mm.yr^{-1}}$.
\end{solution}

\begin{solution}{exo:b1:plate-kinematics:2}
$v = \omega R\sin\Delta = 23.35\sin\Delta$: \qty{11.7}{mm.yr^{-1}} at \ang{30},
\qty{20.2}{mm.yr^{-1}} at \ang{60} and \qty{23.35}{mm.yr^{-1}} at \ang{90}. The
spreading is fastest where the boundary is farthest from the pole, up to
\ang{90}: between Eurasia and North America, at the southern end of the
boundary, near the Azores, \ang{78} from the pole.
\end{solution}

\begin{solution}{exo:b1:plate-kinematics:3}
The subchron lasted $1.072 - 0.988 = \qty{0.084}{\Myr}$, and a half rate in
\unit{mm.yr^{-1}} is also a rate in \unit{km\per\Myr}. (a) $50\times0.084 =
\qty{4.2}{km}$ at \qty{50}{mm.yr^{-1}}; $10\times0.084 = \qty{0.84}{km}$ at
\qty{10}{mm.yr^{-1}}. (b) The stripe straddles the axis:
$2\times10\times0.781 = \qty{15.6}{km}$.
\end{solution}

\begin{solution}{exo:b1:plate-kinematics:4}
The half rates are $66/2.581 = \qty{25.6}{mm.yr^{-1}}$ to the west and
$64/2.581 = \qty{24.8}{mm.yr^{-1}}$ to the east; the full rate is
$130/2.581 = \qty{50.4}{mm.yr^{-1}}$. The
spreading is nearly symmetric: the western flank has grown about
\qty{3}{\percent} faster.
\end{solution}

\begin{solution}{exo:b1:plate-kinematics:5}
$\cos\Delta = \sin\ang{61.8}\sin\ang{85} + \cos\ang{61.8}\cos\ang{85}
\cos(\ang{0} - \ang{139.6}) = 0.8813\times0.9962 + 0.4726\times0.08716\times
(-0.7615) = 0.87795 - 0.03137 = 0.84658$, so $\Delta = \ang{32.16}$ and
$v = 23.35\sin\ang{32.16} = \qty{12.4}{mm.yr^{-1}}$. This is a little above the
measured \qtyrange{11}{12}{mm.yr^{-1}}, but the \omterm{def:b1:plate-kinematics:euler-pole}{Euler vector} is a single best
fit to 453 rates, from which the rates depart by \qty{0.7}{mm.yr^{-1}} on
average, and individual rates near the slow Gakkel Ridge are among the least
precise.
\end{solution}

\begin{solution}{exo:b1:plate-kinematics:6}
(a) $\cos\Delta = \sin\ang{61.8}\sin\ang{39.44} + \cos\ang{61.8}\cos\ang{39.44}
\cos(\ang{169.43}) = 0.20111$: $\Delta = \ang{78.40}$. (b) $v = 23.35
\sin\ang{78.40} = \qty{22.87}{mm.yr^{-1}}$, the \qty{23}{mm.yr^{-1}} measured
just north of the junction. (c) With the formula of the method, the numerator
is $\sin\ang{169.43}\cos\ang{61.8} = 0.08668$ and the denominator
$\cos\ang{39.44}\sin\ang{61.8} - \sin\ang{39.44}\cos\ang{61.8}\cos\ang{169.43}
= 0.97573$: the pole is seen at $\beta = \ang{5.08}$, nearly due north, and
Eurasia moves towards $\ang{95.1}$. (d) The transforms are parallel to the
relative motion: they \omterm{def:b1:field-lab-methods-1:trend-plunge}{trend} about \ang{095}, nearly east--west.
\end{solution}

\begin{solution}{exo:b1:plate-kinematics:7}
(a) $\vec\omega = \omega(\cos\varphi_p\cos\lambda_p, \cos\varphi_p\sin\lambda_p,
\sin\varphi_p)$: $\vec\omega_{\mathrm{EU\,NA}} = (-0.0756, 0.0643, 0.1851)$ and
$\vec\omega_{\mathrm{NB\,NA}} = (0.0333, 0.0282, 0.2289)$, in
\unit{\degree\per\Myr}. (b) $\vec\omega_{\mathrm{NB\,EU}} =
\vec\omega_{\mathrm{NB\,NA}} - \vec\omega_{\mathrm{EU\,NA}} = (0.1089, -0.0361,
0.0438)$, of length
\[
  \sqrt{0.011859 + 0.001303 + 0.001918} = \qty{0.1228}{\degree\per\Myr}.
\]
Its pole is at \omterm{def:b1:spherical-astronomy:geographic}{latitude}
$\arcsin(0.0438/0.1228) = \ang{20.9}$\,N and \omterm{def:b1:spherical-astronomy:geographic}{longitude}
$\arctan(-0.0361/0.1089) = \ang{-18.3}$, that is \ang{18.3}\,W (the
$x$ component being positive, no \ang{180} is added). (c) The boundary between
Nubia and Eurasia offers few direct measurements: between the Azores and
Gibraltar, only the \omterm{def:b1:spherical-astronomy:horizontal}{azimuths} of the Gloria fault. The two spreading motions
across the Atlantic are measured by hundreds of rates, and their difference
inherits their precision.
\end{solution}

\begin{solution}{exo:b1:plate-kinematics:8}
(a) About $\vec e_z$ by \ang{90}: $\cos\ang{90} = 0$, $\vec e_z\cdot\vec e_x = 0$
and $\vec e_z\times\vec e_x = \vec e_y$, so $\vec e_x$ goes to $\vec e_y$. Then
about $\vec e_x$ by \ang{90}: $\vec e_x\cdot\vec e_y = 0$ and
$\vec e_x\times\vec e_y = \vec e_z$, so $\vec e_y$ goes to $\vec e_z$: the point
ends at the north pole. (b) About $\vec e_x$ first: $\vec e_x$ is on the axis,
$\vec e_x\times\vec e_x = \vec 0$ and $(1 - \cos\ang{90})(\vec e_x\cdot\vec e_x)
\vec e_x = \vec e_x$, so it does not move; then about $\vec e_z$ it goes to
$\vec e_y$, at \ang{0}\,N, \ang{90}\,E. (c) The results differ: \omterm{def:b1:plate-kinematics:finite-rotation}{finite
rotations} do not commute. Successive reconstructions must be composed by
multiplying the rotation matrices in the right order, never by adding angles
or poles.
\end{solution}

\begin{solution}{exo:b1:plate-kinematics:9}
(a) In velocity space put $a$ at the origin, $b$ \qty{30}{mm.yr^{-1}} to the
east and $c$ \qty{40}{mm.yr^{-1}} to the south. (b) $\vec v_{CB} =
\vec v_{CA} - \vec v_{BA} = (0, -40) - (30, 0) = (-30, -40)$ (east, north
components): \qty{50}{mm.yr^{-1}} towards the \omterm{def:b1:spherical-astronomy:horizontal}{azimuth}
$\arctan\bigl((-30)/(-40)\bigr) + \ang{180} = \ang{216.9}$, south-west. (c) The
ridge is perpendicular to this direction: it \omterm{def:b1:field-lab-methods-1:trend-plunge}{trends}
$\ang{216.9} - \ang{90} = \ang{126.9}$, that is north-west--south-east.
\end{solution}

\begin{solution}{exo:b1:plate-kinematics:10}
(a) In Rodrigues' matrix, the identity contributes $3\cos\alpha$, the matrix
$\vec n\,\vec n^{\mathsf T}$ the trace $n_x^2 + n_y^2 + n_z^2 = 1$, times
$1 - \cos\alpha$, and the antisymmetric matrix nothing: $\operatorname{tr}R =
3\cos\alpha + 1 - \cos\alpha = 1 + 2\cos\alpha$. (b) Here $\operatorname{tr}R = 0$,
so $\cos\alpha = -1/2$ and $\alpha = \ang{120}$. The axis satisfies $R\vec n =
\vec n$: $R(x, y, z) = (z, x, y)$, so $x = y = z$ and $\vec n =
(1, 1, 1)/\sqrt3$, at \omterm{def:b1:spherical-astronomy:geographic}{latitude} $\arcsin(1/\sqrt3) = \ang{35.26}$\,N and
\omterm{def:b1:spherical-astronomy:geographic}{longitude} \ang{45}\,E. (c) $R\vec e_x = \vec e_y$, $R\vec e_y = \vec e_z$,
$R\vec e_z = \vec e_x$: the three axes are permuted, as a third of a turn about
the diagonal does. Rodrigues' formula with $\alpha = +\ang{120}$ about $\vec n$
indeed gives $\vec e_x\mapsto -\tfrac12\vec e_x + \tfrac12(1, 1, 1) +
\tfrac{\sqrt3}{2}\cdot\tfrac{1}{\sqrt3}(0, 1, -1) = (0, 1, 0)$, so the rotation
is anticlockwise seen from above \ang{35.26}\,N, \ang{45}\,E.
\end{solution}

\begin{solution}{exo:b1:plate-kinematics:11}
(a) On a transform the relative velocity $\vec\omega\times\vec r$ is along the
fault, so parallel to $\vec t$; it is perpendicular to $\vec\omega$, hence
$\vec t\cdot\vec\omega = 0$. The plane perpendicular to $\vec t$ through the
centre contains $\vec r$ (since $\vec t$ is horizontal) and $\vec\omega$: it
is the plane of the \omterm{def:b1:spherical-astronomy:spherical-triangle}{great circle} through the point perpendicular to the fault,
and the pole lies on that circle. Two such planes meet along the line of
$\vec t_1\times\vec t_2$. (b) $\vec t_1 = (0.5436, 0.8367, -0.0673)$ at the
junction and $\vec t_2 = (0.4942, 0.8639, -0.0968)$ in Iceland, so
$\vec t_1\times\vec t_2 = (-0.02285, 0.01936, 0.05612)$, of length $0.06362$.
Its \omterm{def:b1:spherical-astronomy:geographic}{latitude} is $\arcsin(0.05612/0.06362) = \ang{61.9}$\,N and its \omterm{def:b1:spherical-astronomy:geographic}{longitude}
$\ang{180} - \arctan(0.01936/0.02285) = \ang{139.7}$\,E (negative $x$ component):
the pole of the chapter. (c) Badly: the two \omterm{def:b1:spherical-astronomy:spherical-triangle}{great circles} cross at a small
angle, since the \omterm{def:b1:field-lab-methods-1:trend-plunge}{trends} differ by only \ang{8} over \ang{26} of the boundary.
Changing either \omterm{def:b1:field-lab-methods-1:trend-plunge}{trend} by one degree moves the pole by about \ang{15}. Real
determinations combine many transforms with spreading rates, whose variation
along the boundary fixes the angular distance to the pole.
\end{solution}

\begin{solution}{exo:b1:plate-kinematics:12}
(a) $(94 - 4)/3 = \qty{30}{mm.yr^{-1}}$. (b) The differences one year apart are
$24 - 4 = 20$, $25 - 5 = 20$, $26 - 6 = 20$, $65 - 15 = 50$, $74 - 24 = 50$,
$75 - 25 = 50$, $76 - 26 = 50$, $85 - 65 = 20$ and $94 - 74 = 20$ (mm). Sorted:
five values of 20 and four of 50; the median is \qty{20}{mm.yr^{-1}}. (c) Two
positions one year apart are at the same phase of the seasonal oscillation,
which cancels in their difference. A jump affects only the pairs that straddle
it, here four of the nine; as long as fewer than half the pairs straddle a jump,
the median ignores it, while the end-point estimate in (a) takes the whole jump
as motion and comes out \qty{50}{\percent} too high. The method fails when
jumps are so frequent that most one-year pairs contain one, or when the station
moves unsteadily, for instance after a large earthquake.
\end{solution}

\begin{solution}{pb:b1:plate-kinematics:1}
\textbf{1.} $\omega = 0.210\times(\pi/180)/\num{e6} = \qty{3.665e-9}{rad.yr^{-1}}$;
$\omega R = \num{3.665e-9}\times\qty{6.371e9}{mm} = \qty{23.35}{mm.yr^{-1}}$.

\textbf{2.} $\cos\Delta = \sin\ang{61.8}\sin\ang{64.5} + \cos\ang{61.8}
\cos\ang{64.5}\cos(\ang{-157.6}) = 0.79546 - 0.18809 = 0.60737$, so
$\Delta = \ang{52.60}$.

\textbf{3.} $v = 23.35\sin\ang{52.60} = 23.35\times0.7944 =
\qty{18.55}{mm.yr^{-1}}$.

\textbf{4.} The numerator is $\sin\ang{157.6}\cos\ang{61.8} = 0.18011$ and the
denominator
\[
  \cos\ang{64.5}\sin\ang{61.8} - \sin\ang{64.5}\cos\ang{61.8}\cos\ang{157.6} =
  0.37940 + 0.39436 = 0.77376;
\]
both are positive, so $\beta = \arctan 0.2328 =
\ang{13.1}$. Eurasia moves towards $\ang{13.1} + \ang{90} = \ang{103.1}$.

\textbf{5.} East: $18.55\sin\ang{103.1} = \qty{18.07}{mm.yr^{-1}}$; north:
$18.55\cos\ang{103.1} = \qty{-4.20}{mm.yr^{-1}}$, that is \qty{4.20}{mm.yr^{-1}}
to the south.

\textbf{6.} $\cos\Delta = 0.8851$, $\Delta = \ang{27.74}$, $v = 23.35
\sin\ang{27.74} = \qty{10.9}{mm.yr^{-1}}$. The two plates turn about a pole in
eastern Siberia: the Gakkel Ridge is close to it, and the speed falls with
$\sin\Delta$ although the plates are the same.

\textbf{7.} West: $94/9.779 = \qty{9.61}{mm.yr^{-1}}$; east: $97/9.779 =
\qty{9.92}{mm.yr^{-1}}$.

\textbf{8.} $(94 + 97)/9.779 = 191/9.779 = \qty{19.53}{mm.yr^{-1}}$.

\textbf{9.} From C2An: $54.1/2.581 = \qty{20.96}{mm.yr^{-1}}$; from the base of
C1n: $19.8/0.781 = \qty{25.35}{mm.yr^{-1}}$. They do not agree: the younger the
boundary, the higher the rate. A real change of rate is possible, but a steady
excess in the distances, the same at all ages, would do exactly this.

\textbf{10.} The ratio of the west to the east distances is 0.96 to 0.98: the
eastern flank has grown a few per cent more. The full rate depends only on the
sum of the two distances, so the asymmetry does not change it.

\textbf{11.} $(191 - 5.0)/9.779 = \qty{19.02}{mm.yr^{-1}}$;
$(54.1 - 5.0)/2.581 = \qty{19.02}{mm.yr^{-1}}$; $(19.8 - 5.0)/0.781 =
\qty{18.95}{mm.yr^{-1}}$. The three now agree within
\qty{0.1}{mm.yr^{-1}}. The same \qty{5}{km} is a quarter of the \qty{19.8}{km}
of the youngest pair but less than \qty{3}{\percent} of the \qty{191}{km} of the
oldest; old boundaries give the more reliable rates.

\textbf{12.} The corrected half rate is $186/9.779/2 = \qty{9.51}{mm.yr^{-1}}$,
and $9.51\times0.084 = \qty{0.80}{km}$, a third of the \qty{2.5}{km} by which
each boundary is displaced. One would expect so narrow a stripe to be blurred
away. That it is recognized on the Reykjanes Ridge, while it is absent on the
Gakkel Ridge, which is starved of magma, suggests that an abundant supply of
magma records the field faithfully even in narrow stripes.

\textbf{13.} The corrected rates, 18.95 to \qty{19.02}{mm.yr^{-1}}, agree with
the \qty{19.0}{mm.yr^{-1}} of the plate vector; without the correction they
exceeded it by 3 to 33 per cent.

\textbf{14.} $\cos\Delta = \sin\ang{64.139}\sin\ang{59.018} +
\cos\ang{64.139}\cos\ang{59.018}\cos\ang{27.554} = 0.97054$, so $\Delta =
\ang{13.94}$; the chord is $2R\sin(\Delta/2) = 12742\times\sin\ang{6.97} =
\qty{1546}{km}$.

\textbf{15.} $1000/17.5 = 57$ years.

\textbf{16.} The two directions differ by $\ang{123.1} - \ang{114.8} =
\ang{8.3}$: the horizontal component along the \omterm{def:b1:spherical-astronomy:spherical-triangle}{great circle} is
$18.9\cos\ang{8.3} = \qty{18.70}{mm.yr^{-1}}$, and along the chord
$18.70\cos\ang{6.97} = \qty{18.56}{mm.yr^{-1}}$, about \qty{18.6}{mm.yr^{-1}}.

\textbf{17.} The chord lengthens by \qty{17.5}{mm.yr^{-1}}, \qty{6}{\percent}
less than the \qty{18.6}{mm.yr^{-1}} predicted. Both stations stand far from
the plate boundary, so the difference is not local. Either the plates move
slightly more slowly today than on average over the last three million years,
or one of the two estimates carries a small systematic error; the measurements
alone cannot tell. An \omterm{def:b1:plate-kinematics:euler-pole}{Euler vector} computed from satellite positioning alone
predicts \qty{17.1}{mm.yr^{-1}} for this chord, close to the measured value.

\textbf{18.} East: $21.12 - (-1.07) = \qty{22.19}{mm.yr^{-1}}$; north:
$-6.97 - (-0.97) = \qty{-6.00}{mm.yr^{-1}}$. Speed $\sqrt{22.19^2 + 6.00^2} =
\qty{23.0}{mm.yr^{-1}}$, towards $\ang{180} - \arctan(22.19/6.00) =
\ang{105.1}$.

\textbf{19.} The direction is close to the \ang{103} of the plate vector, but
the speed exceeds its \qty{18.6}{mm.yr^{-1}} by a quarter. Both stations stand
inside the zone where the boundary deforms, and H\"ofn, near the largest ice
cap, moves with the rebound of the \omterm{def:b1:layered-earth:layers}{crust} as the ice thins: up by
\qty{12.9}{mm.yr^{-1}}, and sideways, away from the ice. Stations
inside the boundary zone measure the plate motion plus local deformation.

\textbf{20.} $0.210\times10 = \ang{2.10}$.

\textbf{21.} The point travels on a \omterm{def:b1:spherical-astronomy:spherical-triangle}{small circle} of radius $R\sin\Delta$:
$6371\times0.7944\times(2.10\times\pi/180) = \qty{185.5}{km}$. The corrected
width of \omterm{def:b1:layered-earth:layers}{crust} formed since the young edge of C5n, \qty{186}{km} in
\qty{9.8}{Myr}, is nearly the same: the motion has been steady over ten
million years.

\textbf{22.} Components (east, north): Eurasia
$(22.9\sin\ang{95.1}, 22.9\cos\ang{95.1}) = (22.81, -2.04)$, Nubia
$(19.2\sin\ang{103.8}, 19.2\cos\ang{103.8}) = (18.65, -4.58)$. Eurasia relative
to Nubia: $(4.16, 2.54)$, that is \qty{4.9}{mm.yr^{-1}} towards
$\arctan(4.16/2.54) = \ang{58.6}$, east-north-east.

\textbf{23.} A junction of three ridges is \omterm{def:b1:plate-kinematics:triple-junction}{stable} whatever the velocities. The
velocity of question 22 makes Eurasia and Nubia move apart, east-north-east to
west-south-west: their boundary at the junction must open, slowly; it is the
Terceira Rift, whose east-south-east \omterm{def:b1:field-lab-methods-1:trend-plunge}{trend} makes the opening oblique.

\textbf{24.} Part I: \qty{18.6}{mm.yr^{-1}} in the centre of Iceland
(\qty{19.0}{mm.yr^{-1}} at the profile of Part II). Part II: 18.95 to
\qty{19.02}{mm.yr^{-1}} once corrected. Part III: the Reykjavik--Stavanger
chord lengthens by \qty{17.5}{mm.yr^{-1}}, where the plate vector predicts
\qty{18.6}{mm.yr^{-1}}; the pair inside Iceland gives \qty{23.0}{mm.yr^{-1}},
which is not the plate motion. The plate vector and the stripes agree; the
satellites differ by \qty{6}{\percent}. The stripes average over up to ten
million years, the plate vector over \qty{3.16}{Myr}; the satellites measure
the motion of the last twenty to thirty years.

\textbf{25.} The plate vector gives \qty{18.55}{mm.yr^{-1}} across Iceland, and
the corrected stripes \qty{19.0}{mm.yr^{-1}} beside it: to the nearest
millimetre a year, the \omterm{def:b1:plate-kinematics:spreading-rate}{full spreading rate} across Iceland is
\textbf{\qty[text-series-to-math=true]{19}{mm.yr^{-1}}}.
\end{solution}
